\chapter{The Pass Infrastructure}
\label{ch:pebble-passes}

A pass pipeline is an agreement about which facts exist between edits.
Running two correct transformations in sequence can fail when the second consumes an analysis of the first transformation's input.
The first PebbleIR pipeline avoids clever preservation claims and makes its commit boundary explicit.

\section{One transaction around several phases}
\label{sec:pebble-pass-contract}
The supplied \code{pir\_optimize} verifies its input, copies the function, folds constant operations to a fixed point, marks demanded values, removes unmarked definitions, advances the revision, verifies the candidate, and commits.
Neither pass can expose its intermediate construction state through the public function object.
A failed verification leaves the original available for diagnosis.

Algorithm~\ref{alg:pebble-pipeline} presents the same organization as a driver contract.
The current implementation has a fixed pair of phases rather than a plugin registry.
That is enough infrastructure to enforce the invariants without first designing a registration framework for passes that do not exist.

\begin{algorithm}[htbp]
\caption{Running the initial optimization pipeline.}
\label{alg:pebble-pipeline}
\begin{pseudocode}
\Require candidate function $F$ with a revision that can advance
\Ensure verified optimized function or unchanged input
\State verify $F$ and retain its diagnostic on failure
\State $T\gets F$
\State fold eligible constant operations in $T$ until none change
\State mark values demanded by returns and branch conditions in $T$
\State replace unmarked value records with tombstones
\State increment $T.revision$
\State verify $T$; commit only on success
\end{pseudocode}
\end{algorithm}

The observable function includes the CFG and its possible nontermination.
The pipeline deliberately retains every terminator and block.
It is therefore not an aggressive control-flow optimizer, and a constant branch condition does not automatically remove the untaken region.
That restriction keeps reachability and phi-repair obligations out of the first two passes.

\section{Facts need owners and versions}
\label{sec:pebble-pass-facts}
\code{PirFacts} records predecessor sets, dominator sets, the function pointer, and revision.
A cached table is usable only while all three conditions hold: the object belongs to the same logical lifetime, the pointer matches, and the revision matches.
The lifetime condition is external to the current numeric revision scheme and is documented at initialization and parsing boundaries.
It should become an explicit generation field when clients cannot enforce it reliably.

A pass descriptor for a larger implementation should list required facts, preserved facts, representation preconditions, and its failure policy.
For example, a pure value rewrite may preserve the CFG while invalidating value-range facts and reverse use lists.
A block-splitting pass invalidates predecessor tables and may require dominance repair.
Treat preservation declarations as checked engineering claims; periodically recomputing and comparing supposedly preserved facts is a useful debug mode.

\section{A worked cache invalidation}
\label{sec:pebble-pass-worked}
Suppose the running function has revision 9 after construction and verification returns facts tagged 9.
Optimization commits a candidate tagged 10.
The CFG happens to be identical, but the uniform policy invalidates the old cache because its revision no longer matches.
If optimization fails, the original remains revision 9 and its old facts remain associated with that unchanged object.
If the destination is reparsed into a new lifetime, all old facts are discarded regardless of numeric revision.

The test suite checks revision invalidation on successful optimization and byte-for-byte preservation on rejected replacement.
It also checks that a saturated revision counter rejects optimization without mutation.
These boundary tests are more informative than testing only that the happy-path result has the expected arithmetic value.

\section{Instrument decisions before elapsed time}
\label{sec:pebble-pass-measures}
Count allocated slots, live values, successful folds, CFG edges, and VM instructions separately.
They answer different questions.
The running example keeps seven allocated IDs after optimization but drops from seven live values to five.
Its lowered instruction count drops from eight to six.
Calling this a two-slot memory saving would be wrong because tombstones still occupy their slots.
Calling it a measured speedup would also be wrong: instruction count is not elapsed time.

For a future pass manager, record deterministic pass order and optional before-and-after snapshots.
If a failure appears only after phase three, retaining phase two's verified snapshot reduces the search space immediately.
Do not let instrumentation mutate the representation or alter iteration order; otherwise enabling diagnostics may hide the defect under investigation.

\section{Exercises}
\label{sec:pebble-thirteen-exercises}
\begin{enumerate}
\item \label{ex:pebble-failed-pass} State which revision and analysis facts remain valid after a transaction fails verification. Contrast this with a successful parse into the same C object.
\item \label{ex:pebble-three-counts} Give allocated-slot, live-value, and lowered-instruction counts for the running example before and after optimization.
\end{enumerate}

\section{Further reading}
\label{sec:pebble-thirteen-reading}
Morgan's organization of optimization around reusable analyses provides context for pass interfaces~\cite{morgan-1998-optimizing}.
Safonov emphasizes the importance of making a compiler's trust boundaries explicit~\cite{safonov-2010-trustworthy}.
